\section{Preliminaries}
Let $\theta\in\mathbb R^m$ parameterize a density on $\mathbb R^d$,
\begin{equation}
 q_\theta(z)=\sum_{k=1}^K\pi_k(\theta)q_{k,\theta}(z),
 \qquad \pi_k>0,\qquad \sum_k\pi_k=1,
 \qquad \mathcal J(\theta)=\mathbb E_{q_\theta}[f_\theta(Z)].
 \label{eq:objective}
\end{equation}
We assume positive densities on the interior of a parameter-independent support,
differentiability, and integrable derivatives permitting interchange of
differentiation and integration. For transport identities, boundary flux terms
vanish under integration by parts.
Variance comparisons additionally require finite second moments.
Define responsibilities and fixed-sample derivatives by
\begin{align}
 r_k(z)&=\frac{\pi_kq_{k,\theta}(z)}{q_\theta(z)},
 &s_\theta(z)&=\partial_\theta\log q_\theta(z),
 &a_k&=\nabla_\theta\log\pi_k, \label{eq:responsibilities}\\
 h_\theta(z)&=\nabla_z f_\theta(z),
 &d_\theta(z)&=\partial_\theta f_\theta(z).
 \label{eq:derivatives}
\end{align}
Here $\partial_\theta$ holds $z$ fixed; a total derivative through a sample
transformation is written $\nabla_\theta$.
For component reparameterizations, let
$Z_k=T_{k,\theta}(\varepsilon_k)$, with
$\varepsilon_k\sim\rho$ independent of $\theta$, and define the $d\times m$ matrix
\begin{equation}
 V_k(z)=\left.D_\theta T_{k,\theta}(\varepsilon)
              \right|_{\varepsilon=T_{k,\theta}^{-1}(z)}.
 \label{eq:component-velocity}
\end{equation}
Post-stratification requires differentiable component bijections on the common
support. Stratification only requires differentiable maps and uses
$D_\theta T_{k,\theta}(\varepsilon)$ directly, without an inverse.
Let $\operatorname{sg}$ denote an operation that preserves a value and stops its
derivative.

For a fixed joint model $p(x,z)$, the ELBO is
\begin{equation}
 f_\theta(z)=\log p(x,z)-\log q_\theta(z),\qquad
 \mathcal J(\theta)=\log p(x)-\mathrm{KL}(q_\theta\Vert p(\cdot\mid x)).
 \label{eq:elbo}
\end{equation}
Then $d_\theta(z)=-s_\theta(z)$ and
$\mathbb E_{q_\theta}[s_\theta(Z)]
=\nabla_\theta\int q_\theta(z)\,dz=0$.
Thus STL deletes $d_\theta$ from the estimators below. In automatic
differentiation this means using
$f_{\operatorname{sg}(\theta)}(z)=\log p(x,z)-\log q_{\operatorname{sg}(\theta)}(z)$,
while retaining derivatives with respect to $z$.
For a general parameter-dependent $f_\theta$, deleting $d_\theta$ is not valid.
